\section{Unit Subsidies for Negative Dichotomous Valuations}
\label{sec:main}

This section proves the main theorem: for any number of agents $n$ and every
negative dichotomous instance there is an allocation admitting an envy-free
solution with a subsidy of $0$ or $1$ per agent, at most $n-1$ in total, computable in polynomial time. 

% Beyond Definition~\ref{def:dichcost} nothing is
% assumed of the cost functions --- not additivity, not submodularity, not any
% relation between different agents' cost functions --- and the number of chores
% is unrestricted.

Our argument starts from an envy-free \emph{partial} allocation produced by the algorithm of Tao, Wu, Yu, and Zhou~\cite{tao2025existence}, and then completes this partial allocation. We recall in \Cref{sec:m-terminal} the relevant terminal-state properties of their algorithm, including Theorem~6.1 and the three update rules used in its construction. These ingredients are due to Tao et al.~\cite{tao2025existence}. The completion procedure introduced in \Cref{sec:m-completion}, together with the subsequent lemmas and theorems, constitutes our contribution.


% The argument takes as input an envy-free \emph{partial} allocation produced by the algorithm of Tao, Wu, Yu and Zhou~\cite{tao2025existence}, and completes it. That
% algorithm, its Theorem~6.1, and the three update rules recorded in
% \Cref{sec:m-terminal} are due to~\cite{tao2025existence} and are restated here for
% reference; the completion of \Cref{sec:m-completion} onwards, and every lemma
% and theorem after Lemma~\ref{lem:m-terminal}, are ours.

We work in cost form throughout, so an envy-free solution
(Definition~\ref{def:efsolution}) is the requirement $(\mathcal{A},p)$ such that
\begin{equation}
\label{eq:m-ef}
  \cost_i(\bundle i) - \subsidy_i \;\le\; \cost_i(\bundle j) - \subsidy_j
  \qquad \text{for all } i,j \in \agents ,
\end{equation}
and a partial allocation $X$ is envy-free (Definition~\ref{def:ef}) when
\begin{equation}
\label{eq:m-pef}
  \cost_i(X_i) \;\le\; \cost_i(X_j)
  \qquad \text{for all } i,j \in \agents .
\end{equation}
For $S \subseteq \items$ and $e \in \items \setminus S$ write
$\cost_i(e \mid S) := \cost_i(S \cup \set e) - \cost_i(S) \in \set{0,1}$ for a
marginal, and for a partial allocation $X$ write
$R(X) := \items \setminus \bigcup_{i \in \agents} X_i$ for its \emph{residue}.
 Moreover, as $\cost_i$s are integer valued, for all $S,T\subseteq M$, \(\cost_i(S) < \cost_i(T) \implies \cost_i(S) \le \cost_i(T) - 1 .
 \)

To describe the algorithm of Tao, Wu, Yu, and Zhou~\cite{tao2025existence} and our completion procedure, we use the following standard graph-theoretic notions.

\begin{definition}[Strongly connected component and tail component]
\label{def:scc}
Let $G=(V,E)$ be a finite directed graph. Two vertices \(u,v\in V\) are \emph{strongly connected} if either \(u=v\), or there is a directed path from \(u\) to \(v\) and a directed path from \(v\) to \(u\). Strong connectivity is an equivalence relation on $V$, and its equivalence classes are the \emph{strongly connected components} of $G$.

A strongly connected component $S$ is a \emph{tail component} (or \emph{sink component}) if no arc leaves $S$; that is, there is no $(i,j)\in E$ with $i\in S$ and $j\in V\setminus S$.
\end{definition}

The following standard graph-theoretic fact will be used repeatedly.

\begin{lemma}
\label{lem:tail-scc-exists}
Every finite non-empty directed graph has at least one tail strongly connected component.
\end{lemma}


% Every finite non-empty directed graph has at least one tail strongly connected component. Indeed, contracting each strongly connected component to a single vertex yields the condensation graph, which is a finite directed acyclic graph and therefore has a sink.

% \begin{observation}
% \label{obs:tailscc}
% Every finite non-empty directed graph has at least one tail strongly connected component.
% \end{observation}

% \begin{proof}
% Contracting each strongly connected component to a single vertex yields the condensation graph, which is a finite directed acyclic graph and therefore has a sink.
% \end{proof}
% Superseded by Observation~\ref{obs:m-gap} below, which states the same
% integrality gap and additionally the monotonicity of $\cost_i$ used in
% \Cref{sec:m-completion}.
% Recall that as  $\cost_i$s are integer valued, for all $S,T\subseteq M$, \(
%   \cost_i(S) < \cost_i(T) \implies \cost_i(S) \le \cost_i(T) - 1 .
% \)

% \begin{observation}
% \label{obs:m-gap}
% For every $i \in \agents$ and all $S,T \subseteq \items$,
% \[
%   \cost_i(S) < \cost_i(T) \quad\Longrightarrow\quad \cost_i(S) \le \cost_i(T) - 1 .
% \]
% Moreover $\cost_i(S) \le \cost_i(T)$ whenever $S \subseteq T$.
% \end{observation}

% \begin{proof}
% Telescoping from $\cost_i(\emptyset) = 0$ along any ordering of $S$ writes
% $\cost_i(S)$ as a sum of marginals, each lying in $\set{0,1}$ by
% Definition~\ref{def:dichcost}; hence every $\cost_i$ is integer valued, and a
% strict inequality between two of its values is a gap of at least $1$. For
% monotonicity, enumerate $T \setminus S = \set{f_1,\dots,f_q}$ and telescope:
% $\cost_i(T) - \cost_i(S) = \sum_{u=1}^{q} \cost_i\bigl(f_u \mid S \cup
% \set{f_1,\dots,f_{u-1}}\bigr)$, a sum of terms in $\set{0,1}$ by
% Definition~\ref{def:dichcost}, hence non-negative.
% \end{proof}

% --------------------------------------------------------------------------
\subsection{The terminal state of the partial-allocation algorithm}
\label{sec:m-terminal}

\begin{theorem}[\cite{tao2025existence}, Theorem~6.1]
\label{thm:m-twyz}
For every $n$ and all dichotomous cost functions $\cost_1,\dots,\cost_n$ there
is a polynomial-time algorithm returning a partial allocation $X$ that is
envy-free in the sense of \eqref{eq:m-pef} and satisfies
$\abs{R(X)} \le n-1$.
\end{theorem}

% Superseded by the revised statement of the algorithm immediately below.
% Kept for reference; do not compile both -- they share \label{eq:m-graph}.
% The proof of the results in this paper uses not only the statement of Theorem~\ref{thm:m-twyz} but also the state in which the algorithm of~\cite{tao2025existence} halts, so we record that algorithm. It maintains a partial allocation $X = (X_1,\dots,X_n)$, envy-free
% throughout, together with the \emph{equality graph}
% \begin{equation}
% \label{eq:m-graph}
%   G = (\agents, E),
%   \qquad
%   (i,j) \in E \iff i \ne j \ \text{ and } \ \cost_i(X_i) = \cost_i(X_j),
% \end{equation}
% and repeats the following three rules in the stated order, halting only when
% none applies.
% \begin{enumerate}
%   \item[(R1)] If some $e \in R(X)$ and some $i \in \agents$ satisfy
%         $\cost_i(e \mid X_i) = 0$, assign $e$ to $i$.
%   \item[(R2)] If some $e \in R(X)$ and some arc $(i,j) \in E$ lying on a
%         directed cycle of $G$ satisfy $\cost_i(e \mid X_j) = 0$, rotate the
%         bundles along that cycle and assign $e$ to $i$.
%   \item[(R3)] Otherwise select a strongly connected component $S$ of $G$ with
%         no arc of $E$ leaving $S$; such a component exists by
%         \Cref{obs:tailscc}. If $\abs{R(X)} \ge \abs S$, assign one
%         residual chore to each agent of $S$; otherwise halt.
% \end{enumerate}
% The loop runs while $R(X) \ne \emptyset$, so the only exit with
% $R(X) \ne \emptyset$ is the final clause of (R3).

The proof of our main result uses not only the statement of Theorem~\ref{thm:m-twyz}, but also structural properties of the terminal state of the algorithm of~\cite{tao2025existence}. We therefore recall the relevant part of that algorithm. It maintains an envy-free partial allocation $X=(X_1,\dots,X_n)$ together with the \emph{equality graph}
\begin{equation}
\label{eq:m-graph}
  G=(\agents,E),
  \qquad
  (i,j)\in E
  \iff
  i\ne j \ \text{ and }\ \cost_i(X_i)=\cost_i(X_j).
\end{equation}
While $R(X)\ne\emptyset$, the algorithm considers the following rules in the stated order.
\begin{enumerate}
  \item[(R1)] If there exist $e\in R(X)$ and $i\in\agents$ such that
        $\cost_i(e\mid X_i)=0$, assign $e$ to $i$.

  \item[(R2)] Otherwise, if there exist $e\in R(X)$ and an arc
        $(i,j)\in E$ lying on a directed cycle of $G$ such that
        $\cost_i(e\mid X_j)=0$, rotate the bundles along that cycle
        and assign $e$ to $i$.

 \item[(R3)] Otherwise, select a tail strongly connected component
$S$ of $G$. If $\abs{R(X)}\geq\abs S$, assign distinct residual
chores to the agents of $S$, one chore per agent. Otherwise, halt.
\end{enumerate}
Consequently, if the algorithm terminates with \(R(X)\neq\emptyset\), then neither (R1) nor (R2) is applicable at the terminal state. Moreover, if \(S\) denotes the tail strongly connected component selected in the terminating application of (R3), then
\[
    \abs{R(X)}<\abs S.
\]

For the remainder of this section, let $X$ denote the terminal partial allocation of the algorithm, and write
\[
    R:=R(X), \qquad r:=\abs R.
\]
Unless stated otherwise, let $G=(\agents,E)$ denote the equality graph defined in~\eqref{eq:m-graph} at this terminal state. When \(r\ge1\), let \(S\) be the tail SCC selected at the terminating application of (R3); therefore \(r<|S|\).


\begin{lemma}
\label{lem:m-terminal}
Suppose $r \ge 1$. Then
\begin{enumerate}
  \item[(i)] $\cost_i(e \mid X_i) = 1$ for every $i \in \agents$ and every
        $e \in R$;
  \item[(ii)] $\cost_i(e \mid X_j) = 1$ for every arc $(i,j) \in E$ with
        $i,j \in S$ and every $e \in R$;
  \item[(iii)] $\cost_i(X_i) < \cost_i(X_j)$ for every $i \in S$ and every
        $j \in \agents \setminus S$.
\end{enumerate}
\end{lemma}

\begin{proof}
Since $r\geq 1$, the algorithm terminates with nonempty residue.
Therefore, neither \textnormal{(R1)} nor \textnormal{(R2)} is applicable
at the terminal state.

(i) If $\cost_i(e\mid X_i)=0$ for some $i\in\agents$ and $e\in R$,
then \textnormal{(R1)} would be applicable. Since every marginal cost
belongs to $\set{0,1}$ by Definition~\ref{def:dichcost}, it follows that
\[
    \cost_i(e\mid X_i)=1.
\]

(ii) Let $(i,j)\in E$ with $i,j\in S$. Since $S$ is strongly connected,
there is a directed path from $j$ to $i$. Choosing such a path with no
repeated vertices and appending the arc $(i,j)$ yields a directed cycle
containing $(i,j)$. If $\cost_i(e\mid X_j)=0$ for some $e\in R$, then
\textnormal{(R2)} would be applicable to this arc and chore. Hence
$\cost_i(e\mid X_j)\neq 0$, and therefore
\[
    \cost_i(e\mid X_j)=1.
\]

(iii) Let $i\in S$ and $j\in\agents\setminus S$. Since $S$ is a tail
strongly connected component, no arc leaves $S$; in particular,
$(i,j)\notin E$. By~\eqref{eq:m-graph},
\[
    \cost_i(X_i)\neq \cost_i(X_j).
\]
On the other hand, envy-freeness of $X$, as expressed in
\eqref{eq:m-pef}, gives
\[
    \cost_i(X_i)\leq \cost_i(X_j).
\]
The inequality must therefore be strict.
\end{proof}
% --------------------------------------------------------------------------
\subsection{The completion}
\label{sec:m-completion}
Assume $r\geq 1$ and write
\(
    R=\set{e_1,\dots,e_r}.
\) Since $r<\abs S$, we may choose $r$ distinct agents
\(
    T=\set{t_1,\dots,t_r}\subseteq S
\)
arbitrarily. Define the allocation $\alloc=(A_1,\dots,A_n)$ by
\begin{equation}
\label{eq:m-alloc}
  A_i :=
  \begin{cases}
    X_{t_k}\cup\set{e_k}, & i=t_k \quad (1\leq k\leq r),\\[2pt]
    X_i, & i\in\agents\setminus T.
  \end{cases}
\end{equation}
Since the agents $t_1,\dots,t_r$ are distinct and each chore in $R$
is assigned exactly once, the bundles in $\alloc$ remain pairwise
disjoint and their union is $\items$. Hence, $\alloc$ is a complete
allocation. We refer to the agents in $T$ as the \emph{recipients}.

We next construct a subsidy vector $\subsidy\in\set{0,1}^n$ and show
that $(\alloc,\subsidy)$ is envy-free. The following definition
specifies the set of agents who receive one unit of subsidy.

\begin{definition}[Subsidy set]
\label{def:m-P}
Let $P\subseteq\agents$ be the smallest set satisfying
\begin{enumerate}
  \item[(P1)] $T\subseteq P$; and
  \item[(P2)] if $i\in\agents\setminus S$ and $(i,j)\in E$ for some
        $j\in P$, then $i\in P$.
\end{enumerate}
Equivalently, initialise $P:=T$ and repeatedly add any agent
$i\in\agents\setminus S$ for which $(i,j)\in E$ for some $j$ already
in $P$. Continue until no further agents can be added. Since $\agents$
is finite and the set $P$ only grows, this procedure terminates.
Define
\begin{equation}
\label{eq:m-subsidy}
  \subsidy_i :=
  \begin{cases}
    1, & i\in P,\\
    0, & i\in\agents\setminus P.
  \end{cases}
\end{equation}
\end{definition}

The complete procedure is summarised in
Algorithm~\ref{alg:main}. The partial-allocation phase in Line~2 is
the partial-allocation algorithm of Tao et al.~\cite{tao2025existence}; the completion and subsidy
construction in the subsequent lines constitute our contribution.

\begin{algorithm}[H]
\caption{EF1 allocation with at most one unit of subsidy per agent for negative dichotomous valuations}
\label{alg:main}
\begin{algorithmic}[1]
\Require Agents $N=[n]$, chores $M$, and value-oracle access to negative
dichotomous valuations $\{v_i\}_{i\in N}$.
\Ensure A complete allocation $\mathcal{A}$ that is EF1, together with a
subsidy vector $p\in\{0,1\}^n$ such that $(\mathcal{A},p)$ is envy-free and
$\sum_{i\in N}p_i\le n-1$.

\State Set $c_i=-v_i$ for every $i\in N$.
\State Run the partial-allocation algorithm of Tao et al.~\cite{tao2025existence} on $\set{\cost_i}_{i\in\agents}$, retaining the tail strongly connected component selected in the terminating application of  \textnormal{(R3)} if the residue is nonempty.
\State Let $X=(X_1,\dots,X_n)$ be its terminal partial allocation and let
\[
R=\items\setminus\bigcup_{i\in\agents}X_i .
       \]
\If{$R=\emptyset$}
    \State \Return $(X,\mathbf 0)$.
\EndIf

\State Construct the terminal equality graph $G=(\agents,E)$, where
       \[
       (i,j)\in E
       \iff
       i\neq j
       \text{ and }
       \cost_i(X_i)=\cost_i(X_j).
       \]
\State Write $R=\{e_1,\ldots,e_r\}$.
\State Let $S$ be the tail strongly connected component selected in the
       terminating application of \textnormal{(R3)}.
       \Comment{$r<\abs S$}
\State Choose arbitrary distinct agents
      $T=\{t_1,\ldots,t_r\}\subseteq S$.
      \Comment{possible since $r<|S|$}

\For{$k=1,\ldots,r$}
    \State $A_{t_k}\gets X_{t_k}\cup\{e_k\}$.
\EndFor
\For{$i\in N\setminus T$}
    \State $A_i\gets X_i$.
\EndFor

\State Initialise $P\gets T$.
\While{there exist $i\in N\setminus S$ and $j\in P$
       with $(i,j)\in E$ and $i\notin P$}
    \State $P\gets P\cup\{i\}$.
\EndWhile

\For{$i\in N$}
    \State $p_i\gets \mathbf{1}[i\in P]$.
\EndFor

\State $\mathcal{A}\gets(A_1,\ldots,A_n)$.
\State \Return $(\mathcal{A},p)$.
\end{algorithmic}
\end{algorithm}


\begin{lemma}
\label{lem:m-Pstructure}
$P\cap S=T$. Equivalently,
\[P\setminus T\subseteq \agents\setminus S.
\]
Thus, every subsidised agent in $S$ is a recipient.
\end{lemma}

\begin{proof}
Since $T\subseteq P$ by \textnormal{(P1)} and $T\subseteq S$, we have $T\subseteq P\cap S$. Suppose, for contradiction, that there exists $i\in (P\cap S)\setminus T$. Consider
\[
    P':=P\setminus\set{i}.
\]
Since $i\notin T$, the set $P'$ still satisfies \textnormal{(P1)}. We claim that $P'$ also satisfies \textnormal{(P2)}. Let $k\in\agents\setminus S$ and suppose that $(k,j)\in E$ for some $j\in P'$. Since $P'\subseteq P$ and $P$ satisfies \textnormal{(P2)}, we have $k\in P$. Moreover, $k\neq i$, because $k\notin S$ whereas $i\in S$. Hence $k\in P'$.

Thus, $P'$ satisfies both \textnormal{(P1)} and \textnormal{(P2)}, contradicting the minimality of $P$. Therefore $P\cap S=T$.
\end{proof}


% \begin{lemma}
% \label{lem:m-budget}
% $\subsidy\in\set{0,1}^n$ and $\sum_{i\in\agents}\subsidy_i\leq n-1$.
% \end{lemma}

% \begin{proof}
% By \eqref{eq:m-subsidy}, $\subsidy\in\set{0,1}^n$. Moreover, $\abs T=r<\abs S$, so $S\setminus T\neq\emptyset$. Since Lemma~\ref{lem:m-Pstructure} gives $P\cap S=T$, every agent in $S\setminus T$ lies outside $P$ and hence receives zero subsidy. Therefore,
% \[
%     \sum_{i\in\agents}\subsidy_i\leq n-1.
% \]
% \end{proof}






By \eqref{eq:m-subsidy}, $\subsidy\in\set{0,1}^n$. Moreover,
$\abs T=r<\abs S$, so $S\setminus T\neq\emptyset$. Since
Lemma~\ref{lem:m-Pstructure} gives $P\cap S=T$, every agent in
$S\setminus T$ lies outside $P$ and hence receives zero subsidy.
Therefore,
\begin{equation}\label{eq_1}
  \sum_{i\in\agents}\subsidy_i\leq n-1.  
\end{equation}





\begin{lemma}
\label{lem:m-expensive}
Let $i\in S$, $t\in T$, and $e\in R$. Then
\[
    \cost_i(X_t\cup\set e)\geq \cost_i(X_i)+1.
\]
\end{lemma}

\begin{proof}
Since $X$ is envy-free,
\( \cost_i(X_i)\leq \cost_i(X_t).
\)
If the inequality is strict, integrality of the cost functions gives
\( \cost_i(X_t)\geq \cost_i(X_i)+1,
\) and the desired inequality follows from monotonicity of $\cost_i$. Suppose instead that
\(
    \cost_i(X_i)=\cost_i(X_t).
\) If $i=t$, then Lemma~\ref{lem:m-terminal}(i) gives
\[
    \cost_i(e\mid X_i)=1.
\]
If $i\neq t$, then $(i,t)\in E$ by \eqref{eq:m-graph}; since
$i,t\in S$, Lemma~\ref{lem:m-terminal}(ii) gives
\[
    \cost_i(e\mid X_t)=1.
\]
In either case,
\[
    \cost_i(X_t\cup\set e)=\cost_i(X_i)+1.
\]
\end{proof}
% --------------------------------------------------------------------------
\subsection{Envy-freeness of the completion}
\label{sec:m-ef}
\begin{theorem}
\label{thm:m-completion}
Let $X$, $R$, and $S$ be as in \Cref{sec:m-terminal} with $r\geq 1$,
let $\alloc$ be defined by~\eqref{eq:m-alloc}, and let $\subsidy$ be
defined by~\eqref{eq:m-subsidy}. Then $(\alloc,\subsidy)$ is an
envy-free solution.
\end{theorem}

\begin{proof}
Fix $i\in\agents$. Since $X_j\subseteq A_j$ for every $j\in\agents$,
monotonicity of $\cost_i$ and envy-freeness of $X$ give
\begin{equation}
\label{eq:m-mono}
    \cost_i(A_j)
    \geq \cost_i(X_j)
    \geq \cost_i(X_i)
    \qquad\text{for every }j\in\agents.
\end{equation}
We verify
\(
    \cost_i(A_i)-\subsidy_i
    \leq
    \cost_i(A_j)-\subsidy_j
    \)
for every $j\in\agents$.

\smallskip
\noindent\textbf{Case 1: $i\in T$.}
Then $\subsidy_i=1$, and if $e\in R$ is the residual chore assigned
to $i$, Lemma~\ref{lem:m-terminal}(i) gives
\[
    \cost_i(A_i)
    =\cost_i(X_i)+\cost_i(e\mid X_i)
    =\cost_i(X_i)+1.
\]
Hence $\cost_i(A_i)-\subsidy_i=\cost_i(X_i)$. If $j\notin P$, then $j\notin T$, so $A_j=X_j$ and $\subsidy_j=0$. Thus, by envy-freeness of $X$,
\[
    \cost_i(A_j)-\subsidy_j
    =\cost_i(X_j)
    \geq \cost_i(X_i).
\] Suppose instead that $j\in P$. If $j\in T$, then Lemma~\ref{lem:m-expensive} gives
\(
    \cost_i(A_j)\geq \cost_i(X_i)+1.
\) If $j\in P\setminus T$, then
Lemma~\ref{lem:m-Pstructure} implies $j\notin S$. Since $i\in S$, Lemma~\ref{lem:m-terminal}(iii) gives \(\cost_i(X_i)<\cost_i(X_j).\) By integrality, \(
    \cost_i(A_j)=\cost_i(X_j)\geq \cost_i(X_i)+1.
\) Thus, in either case, since $\subsidy_j=1$,
\[
    \cost_i(A_j)-\subsidy_j\geq \cost_i(X_i).
\]

\smallskip
\noindent\textbf{Case 2: $i\in P\setminus T$.}
Here $A_i=X_i$ and $\subsidy_i=1$, so
\(
    \cost_i(A_i)-\subsidy_i=\cost_i(X_i)-1.
\) For every $j\in\agents$, \eqref{eq:m-mono} and
$\subsidy_j\leq 1$ imply
\[
    \cost_i(A_j)-\subsidy_j
    \geq \cost_i(X_i)-1
    =\cost_i(A_i)-\subsidy_i.
\]

\smallskip
\noindent\textbf{Case 3: $i\notin P$.}
Since $T\subseteq P$, we have $i\notin T$, and therefore
$A_i=X_i$ and $\subsidy_i=0$. Hence
\(\cost_i(A_i)-\subsidy_i=\cost_i(X_i).
\) If $j\notin P$, then $A_j=X_j$ and $\subsidy_j=0$, and
envy-freeness of $X$ gives
\[
    \cost_i(A_j)-\subsidy_j
    =\cost_i(X_j)
    \geq\cost_i(X_i)=\cost_i(A_i)-\subsidy_i.
\]
Now suppose that $j\in P$. We claim that
\begin{equation}
\label{eq:m-unpaid-to-paid}
    \cost_i(A_j)\geq \cost_i(X_i)+1.
\end{equation}
If $i\notin S$, envy-freeness of $X$ gives $\cost_i(X_i)\leq\cost_i(X_j)$. Equality would imply $(i,j)\in E$ by~\eqref{eq:m-graph}; since $i\in\agents\setminus S$ and $j\in P$, condition \textnormal{(P2)} would then imply $i\in P$, a contradiction. Hence \(\cost_i(X_i)<\cost_i(X_j), \) and integrality together with monotonicity yield
\[
    \cost_i(A_j)\geq\cost_i(X_j)\geq\cost_i(X_i)+1.
\]
It remains to consider $i\in S$. If $j\in T$, then
\eqref{eq:m-unpaid-to-paid} follows from
Lemma~\ref{lem:m-expensive}. If $j\in P\setminus T$, then
Lemma~\ref{lem:m-Pstructure} gives $j\notin S$, and
Lemma~\ref{lem:m-terminal}(iii), together with integrality, gives
\[
    \cost_i(A_j)=\cost_i(X_j)\geq\cost_i(X_i)+1.
\]
Thus~\eqref{eq:m-unpaid-to-paid} holds in all cases. Since
$\subsidy_j=1$,
\[
    \cost_i(A_j)-\subsidy_j
    \geq\cost_i(X_i)
    =\cost_i(A_i)-\subsidy_i.
\]
The envy-freeness inequality therefore holds for every pair
$i,j\in\agents$.
\end{proof}

\subsection{The EF1 property of the completion}


\begin{corollary}
\label{cor:ef1}
The complete allocation $\mathcal{A}$ constructed in Section~\ref{sec:m-completion} is EF1 in the sense of Definition~\ref{def:ef1chores}, and hence also in the sense of Definition~\ref{def:ef1}. More precisely:
\begin{enumerate}
    \item [(i)] every non-recipient $i\in N\setminus T$ is envy-free in  $\mathcal{A}$; and
    \item [(ii)] for every recipient $t_k\in T$,
    \[
        c_{t_k}(A_{t_k}\setminus\{e_k\})
        \leq
        c_{t_k}(A_j)
        \qquad\text{for every }j\in N.
    \]
\end{enumerate}
\end{corollary}

% \begin{corollary}\label{cor:ef1}
% The complete allocation $\mathcal{A}$ constructed in Section~\ref{sec:m-completion}
% is EF1 in the sense of Definition~\ref{def:ef1chores}, hence also in the
% sense of Definition~\ref{def:ef1}. More precisely, every non-recipient $i\in N\setminus T$ is envy-free in
% $\mathcal{A}$, while for every recipient $t_k\in T$, removing the residual
% chore $e_k$ assigned to $t_k$ eliminates all of $t_k$'s envy.
% \end{corollary}

\begin{proof}
Recall that the partial allocation $X=(X_1,\ldots,X_n)$ is envy-free and that $X_j\subseteq A_j$ for every $j\in N$. Let first $i\in N\setminus T$. Then $A_i=X_i$. Hence, for every $j\in N$, envy-freeness of $X$ and monotonicity of $c_i$ give
\[
    c_i(A_i)
    =
    c_i(X_i)
    \leq
    c_i(X_j)
    \leq
    c_i(A_j).
\]
Thus, every non-recipient is envy-free in $\mathcal{A}$. Now let $i=t_k\in T$. By construction,
\[
    A_i=X_i\cup\{e_k\},
    \qquad\text{and hence}\qquad
    A_i\setminus\{e_k\}=X_i.
\]
Therefore, for every $j\in N$,
\[
    c_i(A_i\setminus\{e_k\})
    =
    c_i(X_i)
    \leq
    c_i(X_j)
    \leq
    c_i(A_j),
\]
where the two inequalities follow from envy-freeness of $X$ and monotonicity of $c_i$, respectively. Thus, removing the residual chore assigned to a recipient eliminates all of her envy. Consequently, $\mathcal{A}$ is EF1.
\end{proof}% --------------------------------------------------------------------------
\subsection{The main theorem}
\label{sec:m-main}

% \begin{theorem}
% \label{thm:m-main}
% For every $n$ and every instance $\langle \agents, \items,
% \set{v_i}_{i \in \agents} \rangle$ with negative dichotomous valuations there
% exist a complete allocation $\alloc$ and a subsidy vector
% $\subsidy \in \set{0,1}^n$ with $\sum_{i \in \agents} \subsidy_i \le n-1$ such
% that $(\alloc,\subsidy)$ is an envy-free solution. Moreover, the allocation $\mathcal{A}$ is EF1. Such an allocation and subsidy vector can be computed in polynomial time given value-oracle access.
% \end{theorem}

% \begin{proof}
% Pass to cost form, $\cost_i = -v_i$, which is dichotomous in the sense of
% Definition~\ref{def:dichcost}. Run the partial-allocation algorithm of Tao et al.~\cite{tao2025existence} and retain
% its terminal partial allocation $X$, the residue $R$ with $r = \abs R$. If $r\geq 1$, construct the terminal equality graph $G$ and fix an arbitrary tail strongly connected component $S$ of $G$, as in
% \Cref{sec:m-terminal}.

% If $r = 0$ then $X$ is a complete allocation satisfying \eqref{eq:m-pef}, so
% $(X, \mathbf 0)$ is an envy-free solution with $\sum_i \subsidy_i = 0 \le n-1$. Since every envy-free allocation is EF1, $X$ is also EF1.

% If $r \ge 1$, construct $T$, $\alloc$ and $\subsidy$ as in
% \Cref{sec:m-completion}. Then $\alloc$ is a complete allocation,
% Theorem~\ref{thm:m-completion} shows $(\alloc,\subsidy)$ is an envy-free
% solution, and (\ref{eq_1}) gives $\subsidy \in \set{0,1}^n$ with
% $\sum_i \subsidy_i \le n-1$. The allocation $\mathcal{A}$ is EF1 by Corollary~\ref{cor:ef1}.

% For the running time, if $r\geq 1$, a tail strongly connected component $S$ of the terminal equality graph can be found in polynomial time by computing its strongly connected components and selecting one with no outgoing arc. Constructing $\alloc$ from \eqref{eq:m-alloc}
% takes $r \le n-1$ assignments. The graph $G$ of \eqref{eq:m-graph} has $n$
% vertices and is determined by the $n^2$ values $\cost_i(X_j)$, one oracle call
% each. Computing $P$ is a backward reachability search in $G$ from $T$
% restricted to vertices outside $S$, which takes $O(n^2)$ time. Hence the work
% after Theorem~\ref{thm:m-twyz} is polynomial, and so is the whole procedure.
% \end{proof}

% As discussed earlier, the bound of Theorem~\ref{thm:m-main} is tight: there are instances on which every complete allocation requires total subsidy
% at least $n-1$.



\begin{theorem}
\label{thm:m-main}
For every instance
$\langle \agents,\items,\set{v_i}_{i\in\agents}\rangle$
with negative dichotomous valuations, there exist a complete allocation
$\alloc$ and a subsidy vector $\subsidy\in\set{0,1}^n$ such that
\[
    \sum_{i\in\agents}\subsidy_i\leq n-1
\]
and $(\alloc,\subsidy)$ is an envy-free solution. Moreover, $\alloc$ is
EF1 even before subsidies are paid. Such an allocation and subsidy vector
can be computed in polynomial time given value-oracle access.
\end{theorem}

\begin{proof}
Pass to the cost representation $\cost_i=-v_i$, which is dichotomous in
the sense of Definition~\ref{def:dichcost}. Run the partial-allocation algorithm of Tao et al.~\cite{tao2025existence}, and let $X$ be its terminal partial
allocation, with residue
\[
    R:=R(X), \qquad r:=\abs R.
\]

If $r=0$, then $X$ is complete and satisfies~\eqref{eq:m-pef}.
Hence $(X,\mathbf 0)$ is an envy-free solution. Since every envy-free
allocation is EF1, $X$ is also EF1.

If $r\geq 1$, let $G=(\agents,E)$ be the terminal equality graph and let $S$ be the tail strongly connected component of $G$ fixed in \Cref{sec:m-terminal}. Construct $T$, $\alloc$, and $\subsidy$ as in \Cref{sec:m-completion}. Then $\alloc$ is complete, Theorem~\ref{thm:m-completion} shows that $(\alloc,\subsidy)$ is envy-free, \eqref{eq:m-subsidy} and \eqref{eq_1} give
\[
    \subsidy\in\set{0,1}^n
    \qquad\text{and}\qquad
    \sum_{i\in\agents}\subsidy_i\leq n-1,
\]
and Corollary~\ref{cor:ef1} shows that $\alloc$ is EF1.

It remains to verify polynomial-time computability. The partial-allocation algorithm of Tao et al.~\cite{tao2025existence} runs in polynomial time. If $r\geq 1$, the
terminal equality graph has $n$ vertices and can be constructed from
the $n^2$ values $\cost_i(X_j)$, requiring one value-oracle query for
each pair $(i,j)$. Its strongly connected components can be computed
in polynomial time, after which a tail component $S$ is obtained by
selecting one with no outgoing arc. By Theorem~\ref{thm:m-twyz},
$r\leq n-1$, so constructing $\alloc$ requires at most $n-1$ additional
assignments. Finally, $P$ can be computed by a backward reachability
search from $T$, restricted to vertices outside $S$, in $O(n^2)$ time. Hence, the entire procedure runs in polynomial time.
\end{proof}

As noted in the introduction, the total-subsidy bound in Theorem~\ref{thm:m-main} is tight: there are instances for which every complete envy-free solution requires total subsidy at least \(n-1\).



\subsection{Pareto optimality and subsidies}

In this section, we show that the unit-subsidy guarantee need not be compatible with Pareto optimality of the underlying allocation. Pareto optimality is evaluated with respect to the allocation of chores, without taking subsidies into account. Formally, a complete allocation $\allocB$ \emph{Pareto dominates} a complete allocation $\alloc$ if
\[
    \cost_i(\allocB_i)\leq \cost_i(\bundle i)
    \qquad\text{for every }i\in\agents,
\]
with strict inequality for at least one agent. A complete allocation is \emph{Pareto optimal} if it is not Pareto dominated by any other complete allocation.

\begin{proposition}
\label{prop_1}
There exist negative dichotomous instances, even with two agents and identical binary submodular cost functions, for which no Pareto-optimal allocation admits an envy-free subsidy vector in $\{0,1\}^n$.
\end{proposition}

\begin{proof}
Consider two agents $\agents=\set{1,2}$ and three chores $\items=\set{a,b,c}$. Both agents have the same cost function
\[
    \cost_1(S)=\cost_2(S)=\min\{\abs S,2\}.
\]
Every marginal cost belongs to $\{0,1\}$, so the cost function is dichotomous. Moreover, it is submodular, since $k\mapsto\min\{k,2\}$ is a nondecreasing concave function of cardinality. It is not additive, since
\[
    \cost_i(\set{a,b,c})=2<3
    =\sum_{g\in\set{a,b,c}}\cost_i(\set g).
\]
Every complete allocation has, up to symmetry, either bundle-size
profile $(3,0)$ or $(2,1)$. First, consider the allocation in which agent $1$ receives all three chores. Its cost profile is $(2,0)$. The possible cost profiles of complete allocations are
\[
    (2,0),\qquad (2,1),\qquad (1,2),\qquad (0,2),
\]
and none of the latter three Pareto dominates $(2,0)$. Hence, the concentrated allocation is Pareto optimal. By symmetry, the allocation in which agent $2$ receives all three chores is also Pareto optimal. However, a concentrated allocation cannot be made envy-free with unit subsidies. If agent $1$ receives all three chores, envy-freeness requires
\[
    2-p_1\leq -p_2,
\]
and hence
\[
    p_1-p_2\geq 2.
\]
Thus, no subsidy vector in $\{0,1\}^2$ suffices. The case in which agent $2$ receives all three chores is symmetric.

Now consider a split allocation, with bundle sizes $(2,1)$ up to symmetry. Its cost profile is $(2,1)$, so it can be made envy-free with subsidies $(1,0)$, up to symmetry. Nevertheless, it is not Pareto optimal. Assigning the remaining chore to the agent who already holds two chores changes the cost profile from $(2,1)$ to $(2,0)$: the cost of that agent remains unchanged, while the other agent's cost strictly decreases.

Thus, the Pareto-optimal allocations are precisely the two concentrated allocations, and each requires a subsidy of at least $2$ for the agent receiving all three chores. Conversely, every allocation that admits an envy-free subsidy vector in $\{0,1\}^2$ is Pareto dominated. Therefore, no allocation is simultaneously Pareto optimal and envy-free with subsidies in $\{0,1\}^2$.
\end{proof}

Thus, the unit-subsidy guarantee of Theorem~\ref{thm:m-main} cannot, in general, be strengthened to require Pareto optimality.
% \begin{remark}
% \label{rem:m-scope}
% Theorem~\ref{thm:m-main} uses Theorem~\ref{thm:m-twyz} together with the four
% properties of its halting state recorded in Lemma~\ref{lem:m-terminal}, which
% are read off the rules (R1)--(R3) of~\cite{tao2025existence} restated in
% \Cref{sec:m-terminal}; the statement of Theorem~\ref{thm:m-twyz} alone does not
% supply them. Envy-freeness is verified directly from \eqref{eq:m-ef} in
% Theorem~\ref{thm:m-completion}, so the argument does not invoke
% Theorem~\ref{thm:hs-characterisation} or Theorem~\ref{thm:hs-minsubsidy}, and
% the subsidy exhibited need not be the pointwise-minimal one. Part (a) of
% Lemma~\ref{lem:m-terminal} is used only in Lemma~\ref{lem:m-budget}: the
% conclusion of Theorem~\ref{thm:m-completion} holds whenever $r \le \abs S$.
% \end{remark}
